% Pursuing Stacks (cotes 134-2 to 134-8): archive cote 134-3, site page 5 (PDF page 6), typescript page 5.
% typescript. Draft compared: ps2.tex:197-255.
% AI-transcribed by Claude Opus 5.5 (claude-opus-5-5, adaptive thinking, effort high), Message Batches API, batch ps1 (msgbatch_01Rv5hCaNTtUt4JHx1yqoHPL), 2026-11-12; not yet checked by eye.
\scanpage{134-3/5}
\mspage{5}

\margin{5}
space, namely sets\inserted{(= 0-stacks)}, categories (or 1-stacks), 2-categories etc. Amon\unsure{g}
these structures, those generalizing groupoids among categories, \inserted{namely}
Gr-stacks of various orders n, play a significant role, especially
for the description of homotopy types, but equally for a non-commutati\hand{ve}
``geometric'' interpretation of the cohomology groups $H^i(X,F)$ of arbitra-%
ry dimension (or ``order), of a topos X with coefficients in an abelian
sheaf F. The reflections in (L) therefore were directly aimed at get-%
ting a grasp on a definition of such Gr-stacks, and whereas it seems to
me to have come to a concrete starting point for such a definition, a
similar reflection for defining \inserted{just} stacks \xout{\illegible{1}} rather than Gr-stacks,
is \inserted{still} lacking\unsure{\over{,}{.}} \xout{and} \over{t}{T}his is \inserted{one} among the manyfold things I have in mind whi\over{he}{le}
sitting down on the present reflections.

Thus \uline{n-stacks}, \uline{relativized over a topos to ``n-stacks over X''}, \uline{are}
\uline{viewed primarily as the natural ``coefficients'' in order to do} \inserted{[non commutative]}\uline{(co)homolo-}%
\uline{gical algebra of dimension} \over{$<$}{$\le$} n over X. The ``integration'' of such
coefficients, in much the same spirit as taking objects $R\hand{\Gamma}F_{\hand{\ast}}$ (\inserted{with $R\Gamma$} the deri-%
ved functor of the sections functor \inserted{$\Gamma$}\hand{)} for complexes of abelian sheaves $F_{\hand{\ast}}$
on X, is here merely the trivial operation of taking sections, namely
the``value'' of the n-stack on the final object of X (or of a representa-%
tive site of X, if X is described in terms of a site). The result of
integration is again an n-stack, whose homotopy \xout{\illegible{1}} groups (with a
dimension shift of n) should be viewed as the cohomology invariants
$H^i(X,F_{\hand{\ast}})$ , where $F_{\hand{\ast}}$ now stands for the n-stack rather than for a com-%
plex of abelian sheaves\hand{(5)}.In my letters (two or three) to Larry Breen in
1975, I develop some heuristics along this point of view, with constant
reference to various geometric situations (mainly from algebraic geome-%
try), providing the motivations. The one motivation maybe which was
the strongest, was the realization that the classical Lefschetz theorems
about comparison of homology and homotopy invariants of a projective
variety, and a hyperplane section - once it was reformulated suitably
so as to get rid of non-singularity assumptions, replaced by suitable
assumptions on cohomological ``depth'' - could be viewed as comparison
statements of ``cohomology'' with coefficients in more or less arbitrary
\uline{stacks}\hand{(6)}.This is carried through completely, within the th\hand{e}n existing
conceptual framework restricted to 1-stacks, in the thesis of Mme Ray-%
naud, a beautiful piece of work. There seems to me to be overwhelming
evidence that her results (maybe her method of proof too ?) should gener-%
lize \over{to}{in} the context of non commutative cohomological algebra of arbitra-%
ry dimension, with a suitable \xout{notion}\inserted[typed]{property} of ind-finiteness as the unique
